\documentclass[10pt,a4paper]{amsart}
\usepackage[margin=19mm]{geometry}
\usepackage{amsmath,amssymb,mathtools}
\usepackage[scr=rsfs,cal=euler]{mathalfa}
\usepackage{xfrac}
\usepackage[unicode,hidelinks,bookmarks=false]{hyperref}
\newcommand{\ie}{i.e.\ }
\newcommand{\cf}{cf.\ }
\newcommand{\resp}{respectively\ }
\newcommand{\Subsection}{\subsection}
\providecommand{\nobreakdash}{\nobreak\hbox{-}}
\setlength{\emergencystretch}{2em}
\multlinegap0pt
\usepackage{amssymb}
\usepackage{enumerate}
\newtheorem{definition}{Definition}
\newtheorem{lemma}{Lemma}
\newtheorem{theorem}{Theorem}
\newtheorem{corollary}{Corollary}
\usepackage{cleveref}
\crefname{definition}{Definition}{Definitions}
\crefname{lemma}{Lemma}{Lemmas}
\crefname{theorem}{Theorem}{Theorems}
\crefname{corollary}{Corollary}{Corollarys}
\newcommand{\R}{\mathbb{R}}
\newcommand{\cK}{\mathcal{K}}
\renewcommand{\d}{\textnormal{d}}
\DeclareMathOperator{\supp}{supp}
\title{Regularity estimates for fractional orthotropic $p$-Laplacians of mixed order}
\author{Jamil Chaker, Minhyun Kim}
\date{Source: arXiv:2104.07507v3 (CC BY 4.0)}
\begin{document}
\maketitle

\noindent\emph{Source and licence.} Regularity estimates for fractional orthotropic $p$-Laplacians of mixed order. Version arXiv:2104.07507v3 (CC BY 4.0), distributed by arXiv under the Creative Commons Attribution 4.0 International licence, http://creativecommons.org/licenses/by/4.0/, as recorded in the arXiv licence metadata for this exact version and shown on the abstract page https://arxiv.org/abs/2104.07507v3. Full licence notice: \texttt{licenses/arxiv-2104.07507-CC-BY-4.0.txt}.\par\smallskip

\noindent\textit{Editorial scope.} Counted unit: the article's lemma lem:iteration (source0.tex lines 692--708). The article's complete original proof of it is delivered with it (source0.tex lines 729--813, one \texttt{proof} environment of 85 lines, which the article places away from the statement). No mathematics is written, repaired or re-derived: the statement and the proof are the article's own lines, in the article's own notation, and no retained line is substituted or re-flowed.  Retained uncounted as the source-local context the proof needs: lines 283--292; lines 296--301; lines 384--390; lines 522--532; lines 712--721.  Nothing was omitted from any retained span, so the package carries no omission record.  No retained line contains a citation, so the wrapper carries no bibliography and no bibliography entry.  No retained line contains a figure environment, an image-inclusion command or drawing code, so no image is delivered and the package declares no asset dependency.  Editorial formatting only, in addition: an AMS \texttt{amsart} preamble that restates the article's own declarations and macros used by the retained lines, namely the environments \texttt{definition}, \texttt{lemma}, \texttt{theorem}, \texttt{corollary} on separate counters, together with the standard packages those lines need.  The retained environments print in this document as Definition 1, Lemma 1, Theorem 1, Corollary 1 in order of appearance, whereas the article prints the same items with the numbering of its own full document.  The complete native source of the article as distributed in the arXiv e-print for this exact version is bundled unchanged as \texttt{sources/110359/source0.tex}.
\begin{definition} \label{def:cut-off}
We say that $(\tau_{x_0,r,\lambda})_{x_0,r,\lambda} \subset C^{0,1}(\R^n)$ is an admissible family of cut-off functions if there is 
$c \ge 1$ such that for all $x_0 \in \R^n$, $r \in (0,1]$ and $\lambda \in (1,2]$, it holds that
\[ \begin{cases}
\supp(\tau) \subset M_{\lambda r}(x_0), \\
\| \tau \|_{\infty} \le 1, \\
\tau \equiv 1 \text{ on } M_r(x_0), \\
\| \partial_k \tau \|_{\infty} \le c \left( \lambda^{s_{\max}/s_k} -1 \right)^{-1}r^{-s_{\max}/s_k} \text{ for every } k \in \lbrace 1 \dots n \rbrace.
\end{cases} \]
\end{definition}

\begin{lemma}\label{lemma:cutoff}
Let $p> 1$, $\Lambda \ge 1$, and $s_1,\dots,s_n\in[s_0,1)$ be given for some $s_0\in(0,1)$. There is $C = C(n,p,\Lambda) > 0$ such that for each $\mu\in\cK(p,s_0,\Lambda)$, every $x_0\in\R^n$, $r \in (0,1]$, $\lambda \in (1,2]$ and every admissible cut-off function $\tau$, the following is true:
\begin{equation*}
\sup_{x\in\R^n} \int_{\mathbb{R}^n} |\tau(y)-\tau(x)|^p \mu(x,\mathrm{d}y) \leq C \left( \sum_{k=1}^n \left( \lambda^{s_{\max}/s_k}-1 \right)^{-ps_k} \right) r^{-ps_{\max}}.
\end{equation*}
\end{lemma}

\begin{theorem}\label{thm:sobolev}
Let $s_1,\dots,s_n\in[s_0,1)$ be given for some $s_0\in(0,1)$. Suppose that $1 < p < n/\bar{s}$ and let $p_{\star} = np/(n-p\bar{s})$. Then, there is a constant $C = C(n, p, p_{\star}, s_0) > 0$ such that
for every $u \in V^{p,\mu_\mathrm{axes}}(\mathbb{R}^n | \mathbb{R}^n)$
\begin{equation} \label{eq:sobolev}
\|u\|_{L^{p_{\star}}(\mathbb{R}^n)}^p \leq C \int_{\mathbb{R}^n} \int_{\mathbb{R}^n} |u(x)-u(y)|^p \mu_{\mathrm{axes}}(x, \mathrm{d}y) \mathrm{d}x.
\end{equation}
\end{theorem}

\begin{corollary} \label{thm:sobolev_loc}
Let $\Lambda \ge 1$ and $s_1,\dots,s_n\in[s_0,1)$ be given for some $s_0\in(0,1)$. Suppose that $1 < p < n/\bar{s}$. There is $C = C(n,p,p_{\star},s_0,\Lambda) > 0$ such that for each $\mu \in \cK(p,s_0,\Lambda)$ and every $x_0 \in \R^n$, $r \in (0,1]$, $\lambda \in (1,2]$ and $u \in H^{p,\mu}_{{M_{\lambda r}(x_0)}}(\R^n)$ it holds
\begin{equation*}
\begin{split}
\|u\|_{L^{{p_{\star}}}(M_r(x_0))}^p\
\leq& C \int_{M_{\lambda r}(x_0)} \int_{M_{\lambda r}(x_0)} |u(x)-u(y)|^p \mu(x, \mathrm{d}y) \mathrm{d}x \\
&+ C \left( \sum_{k=1}^n \left( \lambda^{s_{\max}/s_k}-1 \right)^{-s_k p} \right) r^{-ps_{\max}} \|u\|_{L^p(M_{\lambda r}(x_0))}^p,
\end{split}
\end{equation*}
where ${p_{\star}}$ is defined as in \Cref{thm:sobolev}.
\end{corollary}

\begin{lemma} \label{lem:alg_ineq}
Let $a, b > 0$, $\tau_1, \tau_2 \in [0,1]$, and $t > p-1 > 0$. Then,
\begin{equation*}
\begin{split}
&|b-a|^{p-2}(b-a)(\tau_1^p a^{-t} - \tau_2^p b^{-t}) \\
&\geq c_1 \left| \tau_1 a^{\frac{-t+p-1}{p}} - \tau_2 b^{\frac{-t+p-1}{p}} \right|^p - c_2 |\tau_1-\tau_2|^p \left( a^{-t+p-1} + b^{-t+p-1} \right),
\end{split}
\end{equation*}
where $c_i = c_i(p, t) > 0$, $i=1, 2$, is bounded when $t$ is bounded away from $p-1$.
\end{lemma}

\begin{lemma} \label{lem:iteration}
Let $\Lambda \ge 1$ and $s_1,\dots,s_n\in[s_0,1)$ be given for some $s_0\in(0,1)$. Let $1<p<  n/\bar{s},$ $x_0 \in \R^n$, $r\in(0,1]$, and $\lambda \in (1,2]$. Set $M_r = M_r(x_0)$ and assume $f\in L^{q/(p\bar{s})}(M_{\lambda r})$ for some $q>n$. 
For each $\mu \in \mathcal{K}(p,s_0, \Lambda)$ and every $u \in V^{p,\mu}(M_{\lambda r} | \R^n)$ that satisfies 
\begin{align*}
&\mathcal{E}^\mu(u, \varphi) \geq (f, \varphi) \quad \text{for any nonnegative}~ \varphi \in H^{p,\mu}_{M_{\lambda r}}(\R^n), \\
&u(x) \geq \epsilon \quad \text{a.e. in}~ M_{\lambda r},
\end{align*}
for
\begin{equation*}
\epsilon > r^{\delta}\|f\|_{L^{q/(p\bar{s})}(M_{\lambda r/4})}^{\frac{1}{p-1}} + \left( r^{ps_{\max}} \sup_{x \in M_{(\lambda+1)r/2}} \int_{\mathbb{R}^n \setminus M_{\lambda r}} u_-^{p-1}(y) \mu(x, \d y) \right)^{\frac{1}{p-1}},
\end{equation*}
the following is true for any $t > p-1$,
\begin{equation*}
\left\|u^{-1}\right\|_{L^{(t-p+1)\gamma}(M_r)}^{t-p+1} \leq C \left( \sum_{k=1}^n \left( \lambda^{s_{\max}/s_k}-1 \right)^{-s_k p} \right) r^{-s_{\max}p} \left\|u^{-1}\right\|_{L^{t-p+1}(M_{\lambda r})}^{t-p+1},
\end{equation*}
where $\delta = \frac{ps_{\max}}{p-1} \frac{q-n}{q}$, $\gamma = n/(n-p\bar{s})$, and $C = C(n, p, p_{\star}, q, t, s_0, \Lambda) > 0$ is a constant that is bounded when $t$ is bounded away from $p-1$.
\end{lemma}

\begin{proof} [Proof of \Cref{lem:iteration}]
Let $\tau$ be an admissible cut-off function in the sense of \Cref{def:cut-off}. Since $\tau = 0$ outside $M_{\lambda r}$, the function $\varphi = -\tau^p u^{-t}$ is well defined. Using \Cref{lem:alg_ineq}, we have
\begin{equation*}
\begin{split}
&(f, -\tau^p u^{-t}) \geq \mathcal{E}(u, -\tau^p u^{-t}) \\
&= \int_{M_{\lambda r}} \int_{M_{\lambda r}} |u(y)-u(x)|^{p-2} (u(y)-u(x)) \left( \tau^p(x) u^{-t}(x) - \tau^p(y) u^{-t}(y) \right) \mu(x, \d y)\, \d x \\
&\quad + 2 \int_{M_{\lambda r}} \int_{\mathbb{R}^n \setminus M_{\lambda r}} |u(y)-u(x)|^{p-2} (u(y)-u(x)) \tau^p(x) u^{-t}(x) \mu(x, \d y)\, \d x \\
&\geq c_1 \int_{M_{\lambda r}} \int_{M_{\lambda r}} \left| \tau(x) u^{\frac{-t+p-1}{p}} (x) - \tau(y) u^{\frac{-t+p-1}{p}}(y) \right|^p \mu(x, \d y)\, \d x \\
&\quad - c_2 \int_{M_{\lambda r}} \int_{M_{\lambda r}} |\tau(x)-\tau(y)|^p \left( u^{-t+p-1}(x) + u^{-t+p-1}(y) \right) \mu(x, \d y)\, \d x \\
&\quad - 2 \int_{M_{\lambda r}} \int_{\mathbb{R}^n \setminus M_{\lambda r}} (u(x)-u(y))^{p-1} \tau^p(x) u^{-t}(x) {\bf 1}_{\lbrace u(y) \leq u(x) \rbrace} \mu(x, \d y)\, \d x \\
&=: c_1 I_1 - c_2 I_2 - I_3,
\end{split}
\end{equation*}
where $c_1$ and $c_2$ are constants given in \Cref{lem:alg_ineq}. 
By \Cref{thm:sobolev_loc}, we obtain
\begin{equation*}
\begin{split}
I_1
&= \int_{M_{\lambda r}} \int_{M_{\lambda r}} \left| \tau(x) u^{\frac{-t+p-1}{p}} (x) - \tau(y) u^{\frac{-t+p-1}{p}}(y) \right|^p \mu(x, \d y)\, \d x \\
&\geq C\left\|\tau u^{\frac{-t+p-1}{p}} \right\|_{L^{p_{\star}}(M_{r})}^p -C r^{-ps_{\max}}\left( \sum_{k=1}^n \left( \lambda^{s_{\max}/s_k}-1 \right)^{-s_k p} \right) \left\|\tau u^{\frac{-t+p-1}{p}}\right\|_{L^p(M_{\lambda r})}^p,
\end{split}
\end{equation*}
where
$ p_{\star} = \frac{np}{n-p\bar{s}}$.

For $I_2$, we use \Cref{lemma:cutoff} again to have
\begin{equation*}
\begin{split}
I_2
&= 2\int_{M_{\lambda r}} \int_{M_{\lambda r}} |\tau(x)-\tau(y)|^p u^{-t+p-1}(x) \,\mu(x, \d y)\, \d x \\
&\leq C r^{-ps_{\max}}\left( \sum_{k=1}^n \left( \lambda^{s_{\max}/s_k}-1 \right)^{-s_k p} \right) \left\|u^{-t+p-1}\right\|_{L^1(M_{\lambda r})}.
\end{split}
\end{equation*}
For $I_3$, assuming that $\mathrm{supp}(\tau) \subset M_{(\lambda+1)r/2}$ and using \Cref{lemma:cutoff} we deduce
\begin{equation*}
\begin{split}
I_3
&\leq C \int_{M_{\lambda r}} \int_{\mathbb{R}^n \setminus M_{\lambda r}} \left( u^{p-1}(x)+u_-^{p-1}(y) \right) \tau^p(x) u^{-t}(x) \mu(x, \d y)\, \d x \\
&\leq C r^{-ps_{\max}}\left( \sum_{k=1}^n \left( \lambda^{s_{\max}/s_k}-1 \right)^{-s_k p} \right) \left\|u^{-t+p-1}\right\|_{L^1(M_{\lambda r})} \\
&\quad + C\epsilon^{1-p} \left( \sup_{x \in M_{(\lambda+1)r/2}} \int_{\mathbb{R}^n \setminus M_{\lambda r}} u_-^{p-1}(y) \mu(x, \d y) \right) \left\|u^{-t+p-1}\right\|_{L^1(M_{\lambda r})}.
\end{split}
\end{equation*}
Moreover, we estimate
\begin{equation*}
\begin{split}
|(f, -\tau^p u^{-t})|
&\leq \epsilon^{1-p} \int_{\R^n} |f| \tau^p u^{-t+p-1} \,\d x \\
&\leq \epsilon^{1-p} \|f\|_{L^{q/(p\bar{s})}(M_{\lambda r})} \left\| \tau^p u^{-t+p-1} \right\|_{L^{q/(q-p\bar{s})}(M_{\lambda r})} \\
&= \epsilon^{1-p} \|f\|_{L^{q/(p\bar{s})}(M_{\lambda r})} \left\| \tau u^{\frac{-t+p-1}{p}} \right\|_{L^{pq/(q-p\bar{s})}(M_{\lambda r})}^p.
\end{split}
\end{equation*}
Using Lyapunov's inequality and Young's inequality, we have
\begin{equation*}
\|v\|_{pq/(q-p\bar{s})}^p \leq \|v\|_{p_\star}^{np/q} \|v\|_p^{(qp-np)/q} \leq \frac{n}{q} \omega \|v\|_{p_{\star}}^p + \frac{q-n}{q} \omega^{-n/(q-n)} \|v\|_p^p
\end{equation*}
for any $v \in L^{p_{\star}} \cap L^p$ and any $\omega > 0$. This yields that
\begin{equation*}
\begin{split}
|(f, -\tau^p u^{-t})|
&\leq \epsilon^{1-p} \|f\|_{L^{q/(p\bar{s})}(M_{\lambda r})} \left( \frac{n}{q} \omega \left\| \tau u^{\frac{-t+p-1}{p}} \right\|_{L^{p_\star}}^p + \frac{q-n}{q} \omega^{-n/(q-n)} \left\|\tau u^{\frac{-t+p-1}{p}} \right\|_{L^p}^p \right) \\
&\leq r^{-ps_{\max}\frac{q-n}{q}} \left( \frac{n}{q} \omega \left\| \tau^p u^{-t+p-1} \right\|_{L^{\gamma}} + \frac{q-n}{q} \omega^{-n/(q-n)} \left\|\tau^p u^{-t+p-1} \right\|_{L^1} \right).
\end{split}
\end{equation*}
Combining all the estimates, we have
\begin{equation*}
\begin{split}
&\left\|\tau^p u^{-t+p-1} \right\|_{L^{\gamma}(M_{\lambda r})} \\
&\leq C r^{-ps_{\max}}\left( 1+\sum_{k=1}^n \left( \lambda^{s_{\max}/s_k}-1 \right)^{-s_k p} \right) \left\|u^{-t+p-1}\right\|_{L^1(M_{\lambda r})} \\
&\quad +Cr^{-ps_{\max}\frac{q-n}{q}} \left( \frac{n}{q} \omega \left\| \tau^p u^{-t+p-1} \right\|_{L^{\gamma}(M_{\lambda r})} + \frac{q-n}{q} \omega^{-n/(q-n)} \left\|\tau^p u^{-t+p-1} \right\|_{L^1(M_{\lambda r})} \right).
\end{split}
\end{equation*}
Taking $\omega = \varepsilon_0 r^{-ps_{\max}\frac{q-n}{q}}$ with $\varepsilon_0 > 0$ small enough, we arrive at
\begin{equation*}
\begin{split}
\left\| u^{-1} \right\|_{L^{(t-p+1)\gamma}(M_r)}^{t-p+1}
&\leq \left\|\tau^p u^{-t+p-1} \right\|_{L^{\gamma}(M_{\lambda r})} \\
&\leq C\left( 1+\sum_{k=1}^n \left( \lambda^{s_{\max}/s_k}-1 \right)^{-s_k p} \right) r^{-ps_{\max}} \left\|u^{-1}\right\|_{L^{t-p+1}(M_{\lambda r})}^{t-p+1},
\end{split}
\end{equation*}
where $C$ depends on $n$, $p$, $p_{\star}$, $t$, $s_0$, $q$ and $\Lambda$, and is bounded when $t$ is bounded away from $p-1$. Since $\lambda \leq 2$, we obtain
\begin{equation*}
\sum_{k=1}^n \left( \lambda^{s_{\max}/s_k}-1 \right)^{-s_k p} \geq \sum_{k=1}^n \left( \lambda^{1/s_k}-1 \right)^{-s_k p} \geq \sum_{k=1}^n 2^{-p} = n2^{-p},
\end{equation*}
from which the desired result follows.
\end{proof}

\end{document}
